% Printed 458; principal PDF 465; Google PDF 469.
\par\clearpage
\noindent\makebox[0pt][l]{\fontsize{13}{16}\selectfont 458}\hfill{\small\textbf{HÖHENMESSUNGEN.}}\hfill\null\par\medskip
\begin{center}
{\large Tafel III.}\par\medskip
\begingroup\setlength{\tabcolsep}{12pt}\renewcommand{\arraystretch}{1.05}
\begin{tabular}{r|r@{\hspace{12pt}\vrule width .9pt\hspace{12pt}}r|r@{\hspace{12pt}\vrule width .9pt\hspace{12pt}}r|r}
&$+$&&$+$&&$+$\\
1,9&1&2,8&4&3,4&17\\
2,3&1&2,9&5&3,5&22\\
2,4&2&3,0&7&3,6&27\\
2,5&2&3,1&9&3,7&34\\
2,6&3&3,2&11&3,8&43\\
2,7&3&3,3&14&3,9&54\\
2,8&4&3,4&17&&\\
\end{tabular}\endgroup
\par\bigskip\bigskip
{\large Gebrauch der Tafeln.}
\end{center}
\par\noindent\hspace*{\parindent}$t$, $t'$ Temperatur der Luft; $T$, $T'$ Temperatur des Quecksilbers (nach Réaumur); $b$, $b'$ Barometerstand (in \textls[120]{beliebigem} Maass).

\par\noindent\hspace*{\parindent}Man vermindere $\log b$ und $\log b'$ resp. um $10T$, $10T'$ (als Einheiten der $5^{\mathrm{ten}}$ Decimale betrachtet), und ziehe die so corrigirten Logarithmen von einander ab; der Unterschied sei $=u$. Man addire $\log u$ und $A$, nachdem man, wenn man es für nöthig hält, letzteres nach der zweiten Tafel (die, ebenso wie die dritte, Einheiten in der $5^{\mathrm{ten}}$ Decimale gibt) corrigirt hat; die Summe sei $=v$; diese Grösse erhält noch eine kleine Correction aus Tafel III, von der $v$ selbst das Argument ist. Das so corrigirte $v$ ist der Logarithm des Höhenunterschiedes in Metern. Verlangt man denselben in Toisen, so wird zum Logarithmen noch $9{,}71018$ addirt.
\par\medskip
\noindent\hspace*{.5em}\begin{tabular}{l@{\quad}l@{\quad}l@{\quad}l@{\quad}l}
Beispiel.&$t=15^\circ{,}3,$&$T=14{,}9,$&$b=735{,}581\,\mathrm{mm},$&Polhöhe $=45^\circ$.\\
&$t'=3{,}2,$&$T'=7{,}8,$&$b'=537{,}203$&
\end{tabular}
\par\noindent\hspace*{1.5cm}\begingroup\setlength{\tabcolsep}{2pt}\renewcommand{\arraystretch}{1.2}
\begin{tabular}{r@{\ }c@{\ }r@{\quad}c@{\ }r@{\ }l}
$\log b$&$=$&$2{,}86663$&Corr.&$-149$&\\
$\log b'$&$=$&$2{,}73014$&»&$-78$&\\\cline{2-3}\cline{5-5}
&&$0{,}13649$&»&$-71$&\qquad .\\
$u$&$=$&$0{,}13578$&&&\\\cline{2-3}
$\log u$&$=$&$9{,}13284$&&&\\
$A$&$=$&$4{,}28408$&&&\\
Corr.&$=$&$0$&&&\\\cline{2-3}
$v$&$=$&$3{,}41692$&Corr.&$+18$&$.\ .\ .\ .\ 3{,}41710=\log2612{,}8\,\mathrm{m}.$
\end{tabular}\endgroup
\begin{center}\rule{3.5cm}{.4pt}\end{center}

% Unnumbered title [459]; principal PDF 466; Google PDF 470.
\clearpage\null\vspace*{6.7cm}
\begin{center}{\fontsize{24}{30}\selectfont\textls[160]{HELIOTROP.}}\end{center}
\par\vspace*{11cm}\noindent\makebox[.86\linewidth][r]{58\textsuperscript{*}}\par
% Blank reverse [460]; principal PDF 467; Google PDF 471.
\clearpage\null
% Unnumbered opening [461]; principal PDF 468; Google PDF 472.
\clearpage\vspace*{2.7cm}
\begin{center}
\begin{tabular}{c}\hline
\small\textbf{Göttingische gelehrte Anzeigen.\quad 126. Stück.\quad 9. August 1821.\quad S. 1249--1254.}\\\hline
\end{tabular}\par\bigskip
{\large [Über den Heliotrop.]}
\end{center}
\par\noindent\hspace*{\parindent}Den Kennern der höhern Geodäsie sind die Schwierigkeiten bekannt, sich zur Bildung \textls[120]{grosser} Dreiecke recht zweckmässige Zielpunkte zu verschaffen. Hohe Kirchthürme finden sich in manchen Gegenden nicht in dazu schicklichen Lagen, und auch die vorhandenen bieten oft nicht die gewünschte Gelegenheit zur Aufstellung der Instrumente und zum Centriren der gemessenen Winkel dar; auch ist ihr Bau öfters nicht in dem Maasse regelmässig, wie es zur Erreichung der äussersten Schärfe wünschenswerth ist. Besonders gebaute Signalthürme haben, auch abgesehen von dem Aufwand an Geld und Zeit, welchen ihre Erbauung kostet, mit den Kirchthürmen das gemein, dass sie in solchen Fällen, wo sie sich auf nahen dunkeln Hintergrund projiciren, in beträchtlichen Entfernungen schwer zu sehen und zu pointiren sind, und wenn man ihnen eine helle Farbe gibt, nach der verschiedenen Beleuchtung von der Sonne eine veränderliche höchst nachtheilige Phase zeigen. Ja selbst die vollkommensten Signalthürme, geschwärzte, die sich gegen den Himmel projiciren, sind in \textls[120]{sehr} grossen Entfernungen, wenn man zugleich eine von der Sonne beleuchtete und eine im Schatten befindliche Seitenfläche sieht, nicht gänzlich von einer beschwerlichen Phase frei. Die Messungen bei Nacht mit Hülfe \textsc{Argand}scher Lampen sind zwar diesen Fehlern nicht unterworfen, haben aber dagegen, besonders auf schwer zugänglichen Bergen, andere Inconvenienzen, die zu sehr von selbst einleuchten, als dass es nöthig wäre, sie hier zu berühren.

\par\noindent\hspace*{\parindent}Diese Betrachtungen haben den Hrn. Hofrath \textsc{Gauss} veranlasst, für die

% Printed 462; principal PDF 469; Google PDF 473.
\par\clearpage
\noindent\makebox[0pt][l]{\fontsize{13}{16}\selectfont 462}\hfill{\small\textbf{HELIOTROP.}}\hfill\null\par\medskip
\noindent auf allerhöchsten Befehl im Königreich Hannover auszuführende Gradmessung, auf ein neues Hülfsmittel zu denken, welches, wenn auch nur neben dem genannten, mit Vortheil für die Triangulirung im Grossen anzuwenden wäre. Der Erfolg davon hat seine Erwartungen noch weit übertroffen.

\par\noindent\hspace*{\parindent}Eine auf photometrische Gründe gestützte Untersuchung hatte ihm schon früher die Überzeugung gegeben, dass das von einem nur sehr kleinen Planspiegel reflectirte Sonnenlicht auch in den allergrössten Entfernungen, welche nur bei Triangulirungen vorkommen können, noch hinlängliche Kraft haben müsse, um den schönsten Zielpunkt abzugeben. Um diese Idee zu benutzen, kam es darauf an, ein Instrument anzugeben, mit dessen Hülfe das Sonnenlicht überall genau in jede nöthige Richtung gelenkt werden kann. Es war zugleich die Bedingung zu erfüllen, dass ein solches Instrument überall leicht aufgestellt und gehandhabt werden kann, und dass der Mittelpunkt des reflectirenden Spiegels während der Bewegungen, die gemacht werden müssen, um der fortrückenden Sonne gleichsam zu folgen, stets in absoluter Ruhe bleibt. Ein solches Instrument, welches diese \textls[120]{Lenkung des Sonnenlichts} in jede beliebige Richtung aufs vollkommenste und auf die angezeigte Art auszuführen dient, scheint am schicklichsten den Namen eines \textls[120]{Heliotrops} zu führen, zum wenigsten ebenso schicklich wie zwei bekannte Producte des Pflanzen- und Mineralreichs.

\par\noindent\hspace*{\parindent}Das weitere Nachdenken über diesen Gegenstand hat den Hrn. Hofrath \textsc{Gauss} auf zwei ganz verschiedene Einrichtungen eines Heliotrops geführt; nach der einen ist ein solches Instrument von unserm geschickten Hrn. Inspector \textsc{Rumpf} bereits vortrefflich ausgeführt, und an einem zweiten nach der andern Einrichtung wird von demselben Künstler jetzt gearbeitet. Eine vollständige Beschreibung, die für unsere Blätter sich nicht eignen würde, wird an einem andern Ort gegeben werden. --- Um, noch ehe der erwähnte Heliotrop vollendet war, wirkliche Versuche über die Kraft des reflectirten Sonnenlichts anstellen zu können, kam Hr. Hofrath \textsc{Gauss} noch auf eine dritte Idee, vermittelst welcher jeder Spiegelsextant zu einem ziemlich vollkommenen Heliotrop eingerichtet werden kann, ja allenfalls ohne allen Zusatz, wenn er nur auf ein gutes Stativ gesetzt werden kann, als Heliotrop zu gebrauchen ist. Hr. Hofrath \textsc{Gauss} liess nemlich durch Hrn. Inspector \textsc{Rumpf} an dem grossen Spiegel eines Sextanten einen dritten Planspiegel so befestigen, dass dessen

% Printed 463; principal PDF 470; Google PDF 474.
\par\clearpage
\noindent\null\hfill{\small\textbf{HELIOTROP.}}\hfill\makebox[0pt][r]{\fontsize{13}{16}\selectfont 463}\par\medskip
\noindent Ebene auf der Ebene des Sextanten senkrecht ist, und mit der Ebene des grossen Spiegels einen Winkel macht, der dem Complement des Winkels der Gesichtslinie (die nöthigenfalls erst durch eingezogene feine Kreuzfäden zu bilden ist) gegen die Ebene des kleinen Spiegels zum rechten Winkel gleich ist. Sobald ein solcher Sextant in eine solche Lage gebracht ist, als wollte man die Distanz eines Objects vom Mittelpunkt (ja, insofern alles gut gearbeitet und berichtigt ist, nur von irgend einem Punkt) der Sonnenscheibe messen, gleichviel, welches von beiden direct gesehen wird, reflectirt jener dritte Spiegel das Sonnenlicht nach dem Objecte zu. Steht ein solcher Sextant auf einem guten Stativ, so ist es einer etwas geübten Hand nicht schwer, das reflectirte Sonnenlicht ununterbrochen nach dem gewünschten Punkte hin zu senden. Ist der dritte Spiegel nicht vorhanden, so kann der grosse Spiegel selbst seine Stelle vertreten, wenn man bei völlig unverrückter Ebene die Alhidade um den vorhin erwähnten Winkel schnell vorwärts schiebt (oder nominell auf dem Gradbogen um den doppelten Winkel). Bei dieser letzten Art ist offenbar nur eine unterbrochene Reflexion zu bewirken; doch kann eine geübte Hand, bei Anwendung der nöthigen Sorgfalt und einiger kleiner Kunstgriffe, die hier anzuführen zu weitläuftig sein würde, die Reflexion des Sonnenlichtes nach dem vorgeschriebenen Punkte wohl jedesmal zwei Minuten und darüber anhaltend machen. Beide letztere Arten haben übrigens offenbar die kleine Unvollkommenheit, dass, insofern das Stativ fest steht, der Mittelpunkt des reflectirenden Spiegels nicht in absoluter Ruhe bleibt. In den meisten Fällen wird jedoch dies fast von gar keiner Erheblichkeit sein, so wie man, wenn man es für nöthig hält, auch leicht fortwährend etwas nachhelfen oder davon Rechnung tragen könnte. Auch ist die Anwendbarkeit davon natürlich auf die Winkelentfernung der Sonne vom Object beschränkt, welche die Grösse des Gradbogens des Sextanten vorschreibt.

\par\noindent\hspace*{\parindent}Ehe wir den Erfolg der Versuche, die mit dem Heliotrop angestellt sind, hier anführen, bemerken wir, dass alle zur Reflexion angewandten Spiegel eine Breite von 2 Zoll und eine Höhe von $1\frac14$ Zoll haben. Die Erfahrung hat bestätigt, was Hr. Hofrath \textsc{Gauss} schon aus photometrischen Gründen vorausberechnet hatte, dass bei nur einigermaassen günstigen Umständen grössere Dimensionen ganz unnöthig sein würden, wenigstens für den geodätischen Gebrauch. Bei der zweiten oben erwähnten Einrichtung kann man, übrigens
% Printed 464; principal PDF 471; Google PDF 475.
\par\clearpage
\noindent\makebox[0pt][l]{\fontsize{13}{16}\selectfont 464}\hfill{\small\textbf{HELIOTROP.}}\hfill\null\par\medskip
\noindent auch nach Gefallen und ohne die Dimensionen des Instruments sonst zu vergrössern, einen grössern Spiegel anbringen lassen.
\begin{center}
\begin{tikzpicture}[x=.007cm,y=-.007cm,line cap=round,line join=round]
\path[use as bounding box] (0,0) rectangle (2030,910);
\draw[line width=.65pt] (655,75) circle[radius=.238cm];
\fill (655,75) circle[radius=.0301cm];
\draw[line width=.45pt] (702,125)--(1165,750);
\draw[line width=.45pt] (78,750)--(1165,750)--(1976,740);
\draw[line width=.65pt] (644,735)--(977,735)--(979,740)--(979,760)--(975,765)--(646,766)--(641,760)--(641,741)--cycle;
\draw[line width=.65pt] (608,743)--(641,741)--(642,759)--(610,759)--(606,754)--cycle;
\draw[line width=.5pt] (877,618)--(1413,858);
\draw[line width=.5pt] (1234,607)--(1102,868);
\draw[line width=.5pt,dash pattern=on 0pt off 2.7pt] (884,634)--(1403,867);
\draw[line width=.5pt,dash pattern=on 0pt off 2.7pt] (1239,618)--(1116,867);
\end{tikzpicture}
\end{center}
\par\noindent\hspace*{\parindent}In der Distanz von der hiesigen Sternwarte zum Hohehagen, einem Hauptdreieckspunkte der Gradmessung, (beinahe 2 geographische Meilen) war das Licht vom Heliotrop selbst sowohl, als das von dem zum Heliotrop eingerichteten Sextanten mit blossen Augen, wenn die Sonne hell schien, überaus schön zu sehen; im Fernrohr des Theodolithen war es im Grunde zu stark, und dagegen gab bloss das reflectirte Licht von einer hellen Wolke den schönsten Zielpunkt, der sich denken lässt. Offenbar kann übrigens das reflectirte Sonnenlicht selbst, wo man es wünscht, leicht durch Bedeckung eines Theils des Spiegels nach Gefallen gemässigt werden.

\par\noindent\hspace*{\parindent}In der Entfernung des Hils (eines andern Hauptdreieckspunktes) zum Meridianzeichen der Sternwarte, sehr nahe 5 geographische Meilen, war das Licht beider Instrumente gleichfalls noch mit blossen Augen wie ein schönes Sternchen vortrefflich zu sehen, und bot im Fernrohr des Theodolithen den herrlichsten Zielpunkt dar. Zuweilen bei nebliger Luft, wo von dem Bergrücken des Hils im Fernrohr des Theodolithen ebenso wenig, wie von dem dort erbauten Signalthurm nur eine Spur zu erkennen war, schien das Licht des Heliotrops wie ein prachtvoller Stern im blauen Himmel zu schweben.

\par\noindent\hspace*{\parindent}Die wichtigsten Versuche sind nur erst in den letzten Tagen angestellt. Hr. Professor \textsc{Encke}, Vorsteher der Seeberger Sternwarte, war auf die Einladung des Hrn. Hofrath \textsc{Gauss} hieher gekommen, um den Gebrauch des Spiegelsextanten, ohne dritten Spiegel, als Heliotrop, und die dabei anzuwendenden

% Printed 465; principal PDF 472; Google PDF 476.
\par\clearpage
\noindent\null\hfill{\small\textbf{HELIOTROP.}}\hfill\makebox[0pt][r]{\fontsize{13}{16}\selectfont 465}\par\medskip
\noindent kleinen Kunstgriffe kennen zu lernen, und begab sich sodann auf den Inselsberg, während Hr. Hofrath \textsc{Gauss} die Messungen auf dem Hohehagen anfing. Jener sandte das Sonnenlicht mit dem als Heliotrop gebrauchten Sextanten absatzweise nach dem Hohehagen (Entfernung 85000 Meter oder $11\frac12$ geogr. Meilen), von wo das Sonnenlicht mit dem eigentlichen Heliotrop nach dem Inselsberge gelenkt wurde. Die Versuche und Beobachtungen sind vom 19. bis 29. Julius unter abwechselnd ungünstigen und günstigen Umständen fortgesetzt und haben den allererwünschtesten Erfolg gehabt. Beide Beobachter haben durch das heliotropische Licht die allerschönsten Zielpunkte erhalten, die sich nur irgend denken lassen; häufig erschien es wie ein schönes Sternchen, während man in demselben Fernrohr den Umriss des Berges kaum oder gar nicht wahrnehmen konnte; der eine Beobachter befand sich zuweilen in Nebel und Regen, während das Heliotroplicht von drüben kräftig durchdrang. Ja einige Male glaubten mehrere Anwesende auf dem Hohehagen von vorzüglich scharfer Gesichtskraft das Lichtpünktchen auf dem Inselsberge mit blossen Augen zu erkennen. Wir können noch hinzusetzen, dass die Winkelmessungen selbst, die sich auf das Heliotroplicht bezogen, beiderseitig eine Übereinstimmung gewährt haben, wie sie in einer so grossen Entfernung von keinem andern Signal, es sei denn bei ganz besonders günstigen Umständen, hätte erwartet werden dürfen.

\par\noindent\hspace*{\parindent}Diese Erfahrungen setzen bereits ausser Zweifel, dass bei Anwendung des Heliotroplichts es für die Grösse zu bildender Dreiecke keine Grenzen weiter geben wird, als die die Krümmung der Erde setzt.

\par\noindent\hspace*{\parindent}So wie das Bedürfniss der höhern Geodäsie dieses Instrument veranlasst hat, so beschränken wir uns hier auf Erzählung obiger Erfahrungen, ohne die sich von selbst darbietende Aussicht zu dem künftigen vielleicht noch wichtigern Gebrauch eines den Raum so kräftig durchdringenden Mittels zu telegraphischen Signalisirungen in Krieg und Frieden jetzt weiter zu verfolgen.
\par\vfill\noindent\makebox[0pt][l]{IX.}\makebox[.86\linewidth][r]{59}\par

% Printed 466; principal PDF 473; Google PDF 477.
\clearpage
\noindent\makebox[0pt][l]{\fontsize{13}{16}\selectfont 466}\hfill{\small\textbf{HELIOTROP.}}\hfill\null\par\bigskip\bigskip
\begin{center}
\begin{tabular}{c}\hline
\small\textbf{Astronomisches Jahrbuch für das Jahr 1825.\quad Berlin 1822.\quad S. 103 und 104.}\\\hline
\end{tabular}\par\bigskip
{\large Erfindung eines Heliotrops,\\
vom Hrn. Hofrath Ritter \textsc{Gauss} in Göttingen unterm 26. December 1821\\
eingesandt.}
\end{center}
\par\noindent\hspace*{\parindent}Sie werden mich wegen meines so langen Stillschweigens entschuldigt halten, wenn ich Ihnen sage, dass ich den grössten Theil des Jahres von hier abwesend gewesen bin, und selbst noch im Spätherbst eine Reise nach Altona zur Empfangnahme des \textsc{Ramsden}schen Zenithsectors gemacht habe, von wo ich erst vor kurzem hieher zurückgekommen bin.

\par\noindent\hspace*{\parindent}Bei meiner Triangulation, wo ich bisher an fünf Dreieckspunkten die Winkel gemessen habe, habe ich die Dreiecke so gross wie möglich zu machen gesucht. Über das neue von mir zu diesem Behuf angewandte Hülfsmittel, den \textls[120]{Heliotrop}, und die ersten damit gemachten ins Grosse gehenden Versuche werden Sie die Nachricht in Nr. 126 der hiesigen gelehrten Anzeigen gelesen haben. Seit der Zeit habe ich davon beständig Gebrauch gemacht, nicht allein als Zielpunkt beim Winkelmessen, sondern auch mit nicht weniger glücklichem Erfolg zu telegraphischen Signalisirungen. Die gewaltige Wirkung des reflectirten Sonnenlichts von einem Spiegel von 2 Zoll Breite und $1\frac14$ Zoll Höhe, welches in Entfernungen von 5, 6, $7\frac12$, ja einmal von $9\frac13$ geographischen Meilen mit \textls[120]{blossen Augen} gesehen wurde, pflegt diejenigen, die es zum ersten Male erfahren, und nicht durch theoretische Berechnung darauf vorbereitet sind, gewöhnlich in Erstaunen zu setzen. Bei einem nur einigermaassen günstigen Zustande der Luft gibt es jetzt für die Grösse der Dreiecksseiten keine Grenzen mehr, als die die Krümmung der Erde setzt, zumal wenn man, wie ich es bei zwei neu angefertigten Heliotropen von ganz verschiedener Construction gethan habe, den Spiegeln noch etwas grössere Dimensionen gibt.
\begin{center}\rule{3.5cm}{.4pt}\end{center}

% Printed 467; principal PDF 474; Google PDF 478.
\par\clearpage
\noindent\null\hfill{\small\textbf{BRIEFWECHSEL.\quad HELIOTROP.}}\hfill\makebox[0pt][r]{\fontsize{13}{16}\selectfont 467}\par
\vspace*{4.2cm}
\begin{center}{\fontsize{19}{24}\selectfont\textls[80]{BRIEFWECHSEL.}}\end{center}
\vspace*{1.5cm}
\begin{center}{\large\textsc{Gauss} an \textsc{Olbers}.\quad Göttingen, 1. Julius 1821.}\end{center}
\par\noindent\hspace*{\parindent}.\ .\ .\ .\ .\ .\enspace Es lag mir inzwischen daran, vorerst nur über die Strahlkraft der Spiegel selbst einige Erfahrungen zu erhalten. Ich habe erst mancherlei versucht. Ich befestigte einen Spiegel am Deckel des Fernrohrs eines Theodolithen und suchte durch im voraus mühsam berechnete Azimuthe und Höhen dem Spiegel die richtige Lage zu geben, um das Licht nach einer bestimmten Richtung zu werfen. Dies misslang aber gänzlich. Der Deckel, etwas hart gehend, konnte nicht mit Sicherheit immer wieder in dieselbe Lage gebracht werden, sondern es blieben darin Differenzen von $20'$, die dies Verfahren ganz unbrauchbar machten, wenn nicht der Spiegel auf eine solidere Art am Fernrohr befestigt wurde, so dass dieses offen blieb. Inzwischen brachte mich der Verdruss über die verlorne Mühe auf eine andere Idee, die vollkommen gelungen ist. Der blosse Spiegelsextant auf einem guten Stativ leistet schon das Verlangte, obwohl nicht so vollkommen wie ein eigentlicher Sonnenspiegel. Ist der Winkel, den die Gesichtslinie (die, wenn sie nicht schon vorhanden, erst durch einen Faden oder ein Fadenkreuz dargestellt werden muss) mit dem kleinen Spiegel macht, $=90^\circ-a$, und ist Sonnenbild und Object, wohin das Licht zu werfen, auf gewöhnliche Art zur Berührung gebracht, als wollte man den Winkel messen --- gleichviel ob ersteres oder letzteres direct gesehen ---, so braucht man nur bei unverrückter Ebene die Alhidade um $a$ (oder nominell $2a$) vorzurücken und hat seinen Zweck erreicht. Man kann bei einiger Übung die Stellung leicht so machen, dass jene Coincidenz
\par\medskip\noindent\makebox[.86\linewidth][r]{59\textsuperscript{*}}\par
% Continuation on printed 468.
